\documentclass[11pt,a4paper]{article}
\usepackage[hyperref]{acl}
\usepackage{times}
\usepackage{latexsym}
\usepackage{booktabs}
\usepackage{amsmath}
\usepackage{amssymb}
\usepackage{graphicx}
\usepackage{multirow}
\usepackage{xcolor}
\usepackage{enumitem}
\usepackage[T1]{fontenc}
\usepackage[utf8]{inputenc}
\usepackage{microtype}

\title{Generation Depth Fingerprinting: Classifying Recursive LLM Output\\Beyond Binary Detection}

\author{Anonymous}

\begin{document}
\maketitle

\begin{abstract}
Web-scraped training corpora increasingly contain LLM-generated text of unknown provenance depth---material that has been recursively generated one or more times. Current detectors offer only binary real-vs-synthetic classification, conflating qualitatively different generations. We formulate \textbf{generation depth classification} (gen-0 / gen-1 / gen-2+) and characterize \textbf{differential depth response} (DDR): central-tendency features (entropy mean, TTR, hapax ratio) drift monotonically with depth ($|\rho| > 0.60$), while distributional shape features (entropy max, tail decay) remain near-stationary ($|\rho| < 0.10$), with a ${\sim}10\times$ gap in KS statistics between groups. A lightweight 7-feature ensemble achieves leak-free binary F1 = 0.694 (95\% CI [0.688, 0.700]) under 8-generator leave-one-generator-out (LOGO) evaluation across 5 model families (0.6B--8B), and 0.830 on cross-model chains. DDR is robust across tested decoding strategies (greedy, top-k) and three reference models from different families. We show that depth-aware data curation is statistically equivalent to random subsampling for downstream CLM quality (TOST $p = 1.1 \times 10^{-4}$ at $\pm 2\%$ margin), establishing depth fingerprinting as a \textit{diagnostic} tool for provenance characterization rather than a curative one for data selection. The signal degrades substantially under instruction tuning, restricting current scope to base models.
\end{abstract}

%------------------------------------------------------------
\section{Introduction}
\label{sec:intro}
%------------------------------------------------------------

Large language models (LLMs) are increasingly deployed at scale, and their outputs feed back into the training pipelines of subsequent models. Recent estimates suggest 7--17\% of scientific papers may contain LLM-generated content, and web crawls used for pretraining are progressively contaminated with synthetic text of varying provenance. This recursive contamination---where model-generated text is regenerated by the same or different models---raises fundamental questions about data quality and the distributional properties of recursively generated text.

Current AI-text detection methods frame the problem as binary classification: human-written vs.\ machine-generated. This conflation obscures meaningful differences. First-generation text (gen-1) retains much of the source distribution's lexical properties, while deeper generations (gen-2+) exhibit measurably reduced vocabulary diversity and compressed perplexity distributions. Existing detectors such as FastDetectGPT \cite{bao2024fast}, Binoculars \cite{hans2024binoculars}, and LD-Score cannot distinguish these qualitatively different materials.

We move beyond binary detection and formulate \textbf{generation depth classification}: given a text passage and a reference language model, estimate whether the text is original (gen-0), first-generation synthetic (gen-1), or deeper-generation synthetic (gen-2+). Our central empirical finding is \textbf{differential depth response} (DDR): features measuring central tendency and lexical diversity drift monotonically with generation depth, while features capturing distributional shape remain approximately stationary. This is consistent with recursive generation inducing a location shift in the NLL distribution while preserving its shape.

Our contributions are:
\begin{enumerate}[nosep,leftmargin=*]
\item \textbf{Task formulation}: Generation depth classification as a new problem beyond binary detection, with a rigorous LOGO evaluation protocol that prevents generator leakage.
\item \textbf{Empirical characterization of DDR} across 5 generators from 3 families (0.6B--8B), validated with 3 reference models. DDR is robust across tested decoding strategies (greedy, top-k).
\item \textbf{Negative finding}: Depth-aware curation does not improve downstream CLM quality over random subsampling (TOST-confirmed equivalence), positioning depth fingerprinting as diagnostic rather than curative.
\end{enumerate}

%------------------------------------------------------------
\section{Related Work}
\label{sec:related}
%------------------------------------------------------------

\paragraph{AI-Text Detection.}
Perplexity-based detectors leverage the observation that model-generated text has lower perplexity under the generating model. GLTR \cite{gehrmann2019gltr} visualizes token-level statistics; DetectGPT \cite{mitchell2023detectgpt} uses perturbation-based log-probability curvature; FastDetectGPT \cite{bao2024fast} replaces perturbations with conditional sampling for efficiency; Binoculars \cite{hans2024binoculars} contrasts two models' cross-entropy scores. All produce binary decisions and do not estimate generation depth. When repurposed as ordinal depth predictors (using continuous scores as features), these methods show high cross-generator variance and substantially lower performance than multi-feature approaches.

\paragraph{Model Collapse and Iterative Generation.}
\citet{shumailov2024ai} show that training on model-generated data leads to ``model collapse'' with progressive distribution narrowing. The ``LLM as Broken Telephone'' framework (ACL 2025) demonstrates that paraphrasing chains converge to attractor states within ${\sim}3$ iterations under semantic-preservation constraints. Our work differs in studying non-paraphrastic continuation chains, where the absence of semantic constraints permits monotonic drift beyond the attractor horizon.

\paragraph{Data Quality and OOD Detection.}
Perplexity-based data scoring is standard in pretraining pipelines. Our reference-model features (entropy statistics, tail decay) are related to OOD detection methods that use NLL distribution properties. The distinction is that we exploit the \textit{ordinal} structure of depth-induced distributional shift rather than binary in-vs-out classification.

\paragraph{Authorship Attribution and Stylometry.}
Surface features such as TTR and hapax legomena ratio have long histories in authorship attribution. We repurpose these as depth indicators, showing they capture vocabulary impoverishment across recursive generation.

%------------------------------------------------------------
\section{Method}
\label{sec:method}
%------------------------------------------------------------

\subsection{Problem Formulation}

Given a text passage $x$ and a reference model $M_\text{ref}$, we classify $x$ into one of three depth classes: gen-0 (original human text), gen-1 (first-generation LLM output), or gen-2+ (deeper recursive generations). We define \textbf{generation depth} operationally: depth $d$ text is produced by feeding depth $d{-}1$ text as a prompt to a generator model, starting from human-written seeds at $d{=}0$.

\subsection{Feature Design}

We extract 7 features per passage, divided into two groups motivated by the DDR hypothesis:

\paragraph{Surface features (2):} Type-Token Ratio (TTR) and hapax legomena ratio. These capture vocabulary diversity, which decreases monotonically with generation depth as models converge toward high-frequency token distributions.

\paragraph{Reference features (5):} Given reference model $M_\text{ref}$, we compute token-level negative log-likelihood (NLL) values and extract: entropy mean, entropy standard deviation, entropy minimum, entropy maximum, and tail decay rate. These capture properties of the NLL distribution that respond differentially to generation depth.

\subsection{Differential Depth Response (DDR)}

We classify features into three groups based on their Spearman correlation $\rho$ with depth and KS test statistics:
\begin{itemize}[nosep,leftmargin=*]
\item \textbf{Drifting} ($|\rho| > 0.50$, KS $> 0.50$): hapax ratio ($\rho = -0.63$), TTR ($-0.61$), entropy mean ($-0.60$).
\item \textbf{Moderate} ($0.10 < |\rho| < 0.50$): entropy min ($-0.35$), entropy std ($-0.12$).
\item \textbf{Near-stationary} ($|\rho| < 0.10$, KS $< 0.06$): entropy max ($0.06$), tail decay ($0.04$).
\end{itemize}

The ${\sim}10\times$ gap in KS statistics between drifting features (0.516--0.537) and near-stationary features (0.043--0.052) provides non-parametric confirmation of the three-tier DDR structure.

\subsection{Classifier}

We train an XGBoost classifier on the 7-feature vector. Evaluation follows the leave-one-generator-out (LOGO) protocol: models are trained on chains from $N{-}1$ generators and tested on the held-out generator, ensuring no generator-specific artifacts leak into evaluation. The primary metric is \textbf{leak-free binary F1} (gen-1 vs.\ gen-2+), excluding gen-0 to avoid seed-sharing leakage between gen-0 samples across generators.

\subsection{Chain Construction}

We construct recursive generation chains of depth 0--4 using three prompt types:
\begin{itemize}[nosep,leftmargin=*]
\item \textbf{Continuation}: Direct text continuation from the previous depth's output.
\item \textbf{Elaboration}: Prompted to elaborate on the previous depth's content.
\item \textbf{Inspired rewriting}: Prompted to rewrite inspired by the previous content.
\end{itemize}

Seeds are drawn from multiple domains (Wikipedia, news, academic, code) with 500--1000 seeds per domain.

%------------------------------------------------------------
\section{Experiments}
\label{sec:experiments}
%------------------------------------------------------------

\subsection{Setup}

We evaluate on 5 base-model generators spanning 3 families: Pythia-2.8B, Pythia-1B (Pythia family), OLMo-2-1B (OLMo), Qwen3-1.7B, and Qwen3-0.6B (Qwen). An additional large-scale experiment includes Qwen3-8B. The default reference model is Pythia-410M; cross-reference experiments use OLMo-2-1B and Qwen3-0.6B as alternative references. Total dataset: 18,423 passages across 5 generators.

\subsection{Feature Ablation and DDR Analysis}
\label{sec:ablation}

Table~\ref{tab:ablation} shows the feature ablation results on Pythia-2.8B using 5-fold cross-validation.

\begin{table}[t]
\centering
\small
\begin{tabular}{lcc}
\toprule
\textbf{Feature Set} & \textbf{Macro F1} & \textbf{95\% CI} \\
\midrule
Full (7 features) & 0.820 & [0.807, 0.833] \\
Surface-only (TTR + hapax) & 0.763 & [0.747, 0.774] \\
Reference-only (5 entropy + tail) & 0.719 & [0.706, 0.733] \\
\bottomrule
\end{tabular}
\caption{Feature ablation (5-fold CV, Pythia-2.8B). McNemar's test confirms Full $>$ Surface $>$ Reference (all $p \approx 0$).}
\label{tab:ablation}
\end{table}

The full feature set significantly outperforms both subsets, confirming feature complementarity. Surface features contribute more individually (+4.4pp over reference-only), but reference features provide independent complementary signal (+5.7pp when added to surface).

Table~\ref{tab:ddr} presents the DDR feature classification. Drifting features exhibit KS statistics 0.516--0.537, while near-stationary features show 0.043--0.052---a ${\sim}10\times$ gap that provides non-parametric confirmation of DDR.

\begin{table}[t]
\centering
\small
\begin{tabular}{lccc}
\toprule
\textbf{Feature} & \textbf{Spearman $\rho$} & \textbf{KS} & \textbf{DDR Group} \\
\midrule
hapax\_ratio & $-0.63$ & 0.537 & drifting \\
TTR & $-0.61$ & 0.516 & drifting \\
entropy\_mean & $-0.60$ & 0.528 & drifting \\
entropy\_min & $-0.35$ & 0.052 & moderate \\
entropy\_std & $-0.12$ & 0.043 & moderate \\
entropy\_max & $+0.06$ & 0.048 & near-stationary \\
tail\_decay & $+0.04$ & 0.045 & near-stationary \\
\bottomrule
\end{tabular}
\caption{DDR feature classification. Drifting features show strong monotonic depth correlation and high distributional shift; near-stationary features show negligible response.}
\label{tab:ddr}
\end{table}

Median perplexity decreases monotonically with depth (13.50, 4.96, 3.00, 2.07, 1.72 for depths 0--4), validating the location-shift interpretation: recursive generation progressively moves the NLL distribution toward lower values.

\subsection{Cross-Generator Evaluation}
\label{sec:cross_gen}

Table~\ref{tab:logo} reports per-generator results under 5-generator LOGO evaluation.

\begin{table}[t]
\centering
\small
\begin{tabular}{llccc}
\toprule
\textbf{Generator} & \textbf{Family} & \textbf{Size} & \textbf{Bin.~F1} & \textbf{3-cl.~F1} \\
\midrule
Pythia-2.8B & Pythia & 2.8B & 0.724 & 0.812 \\
Pythia-1B & Pythia & 1B & 0.721 & 0.851 \\
Qwen3-1.7B & Qwen & 1.7B & 0.688 & 0.811 \\
OLMo-2-1B & OLMo & 1.5B & 0.662 & 0.752 \\
Qwen3-0.6B & Qwen & 0.6B & 0.656 & 0.819 \\
\midrule
\textbf{Average} & & & \textbf{0.694} & \textbf{0.804} \\
\bottomrule
\end{tabular}
\caption{LOGO per-generator results. Binary F1 = gen-1 vs gen-2+ (leak-free). Per-generator range: 0.656--0.724 (6.8pp spread).}
\label{tab:logo}
\end{table}

The LOGO binary F1 of 0.694 [0.688, 0.700] demonstrates that the depth signal generalizes to unseen generators. The 6.8pp per-generator spread (0.656--0.724) reflects generator-specific variation in depth signal strength. Feature importance is stable across generators: hapax ratio (0.289), TTR (0.168), entropy mean (0.160).

\paragraph{Cross-family transfer.} Leave-one-family-out (LOFO) evaluation confirms robust cross-family generalization (Table~\ref{tab:lofo}). The average LOFO binary F1 of 0.686 is only 0.8pp below LOGO, indicating that depth features transfer across model families.

\begin{table}[t]
\centering
\small
\begin{tabular}{lcc}
\toprule
\textbf{Held-out Family} & \textbf{Binary F1} & \textbf{$\Delta$ vs LOGO} \\
\midrule
Pythia & 0.678 & $-0.016$ \\
OLMo & 0.691 & $-0.003$ \\
Qwen & 0.689 & $-0.005$ \\
\midrule
\textbf{Average} & \textbf{0.686} & $-0.008$ \\
\bottomrule
\end{tabular}
\caption{LOFO cross-family transfer. Delta vs LOGO is minimal ($-0.8$pp average).}
\label{tab:lofo}
\end{table}

\subsection{Cross-Reference Model Validation}
\label{sec:cross_ref}

To verify DDR is not an artifact of the Pythia-410M reference, we replicate with OLMo-2-1B and Qwen3-0.6B as reference models (Table~\ref{tab:ref_model}).

\begin{table}[t]
\centering
\small
\begin{tabular}{lcccc}
\toprule
\textbf{Reference} & \textbf{Size} & \textbf{3-cl.~F1} & \textbf{ent.~drift} & \textbf{$|$tail $\rho|$} \\
\midrule
Pythia-410M & 410M & 0.803 & [$-1.19$, $-1.40$] & 0.05 \\
OLMo-2-1B & 1.5B & 0.807 & [$-0.61$, $-0.91$] & 0.07 \\
Qwen3-0.6B & 0.6B & 0.816 & [$-0.96$, $-1.35$] & 0.04 \\
\bottomrule
\end{tabular}
\caption{Reference model comparison. DDR direction is consistent across all references; LOGO F1 varies only 1.3pp. Drift magnitude is reference-dependent (OLMo ${\sim}50\%$ of Pythia), but the drifting/stationary classification is preserved.}
\label{tab:ref_model}
\end{table}

All three references yield consistent DDR: entropy drift is uniformly negative, tail decay $|\rho|$ remains below 0.07, and LOGO 3-class F1 ranges 0.803--0.816 (1.3pp spread). Greedy decoding shape stability is confirmed across all references (mean $|d|$ = 0.05--0.06). While drift magnitude varies by reference (OLMo reference drift is ${\sim}50\%$ of Pythia's), the qualitative DDR structure---drifting vs.\ near-stationary---is preserved.

\subsection{Baseline Comparison}
\label{sec:baselines}

Table~\ref{tab:baselines} compares the 7-feature ensemble against existing detectors repurposed as ordinal depth predictors.

\begin{table}[t]
\centering
\small
\begin{tabular}{lcccc}
\toprule
\textbf{Method} & \textbf{F1} & \textbf{95\% CI} & \textbf{$\rho$} & \textbf{Type} \\
\midrule
7-feat ensemble & 0.798 & [.786, .799] & --- & multi \\
Avg-LogProb & 0.765 & [.616, .820] & $-0.79$ & single \\
LD-Score & 0.721 & [.588, .770] & $-0.79$ & single \\
FastDetectGPT & 0.627 & [.511, .762] & $-0.43$ & single \\
Binoculars & 0.579 & [.405, .664] & $-0.20$ & single \\
\bottomrule
\end{tabular}
\caption{Ordinal baseline comparison (LOGO 3-class F1). The 7-feature ensemble outperforms all single detectors (+3.3pp vs best, $p < 0.001$) with substantially narrower CI, reflecting robust cross-generator behavior.}
\label{tab:baselines}
\end{table}

The 7-feature ensemble outperforms the best single detector (Avg-LogProb) by 3.3pp ($p < 0.001$, bootstrap). More importantly, its CI [0.786, 0.799] is far narrower than Avg-LogProb's [0.616, 0.820], demonstrating consistent performance across generators. Single detectors designed for binary classification show high cross-generator variance when repurposed for ordinal depth prediction---their ordinal utility is limited.

\subsection{Downstream Utility: A Negative Finding}
\label{sec:downstream}

We test whether depth-aware data curation improves downstream language modeling quality by fine-tuning OLMo-1B on depth-stratified subsets (3 seeds each). Table~\ref{tab:downstream} reports perplexity results.

\begin{table}[t]
\centering
\small
\begin{tabular}{lccc}
\toprule
\textbf{Variant} & \textbf{$n$} & \textbf{PPL} & \textbf{TOST $p$} \\
\midrule
All depths & 12,299 & 2.335 $\pm$ 0.000 & --- \\
Size-matched random & 3,951 & 2.367 $\pm$ 0.000 & ref \\
Depth 0+1 only & 3,951 & 2.387 $\pm$ 0.001 & $1.1{\times}10^{-4}$ \\
Binary-filtered & 3,951 & 2.389 $\pm$ 0.001 & ${\sim}10^{-4}$ \\
Depth 0 only & 797 & 2.681 $\pm$ 0.001 & reject \\
\bottomrule
\end{tabular}
\caption{Downstream CLM quality by training data composition. TOST at $\pm 2\%$ margin confirms statistical equivalence between size-matched random subsampling and depth-filtered variants. Depth 0 only is worse due to sample size, not depth quality.}
\label{tab:downstream}
\end{table}

Depth-aware curation does not improve downstream CLM quality. At equal sample count ($n = 3{,}951$), random subsampling (PPL = 2.367) is statistically equivalent to depth 0+1 filtering (PPL = 2.387; TOST $p = 1.1 \times 10^{-4}$ at $\pm 2\%$ margin). The depth-0-only variant (PPL = 2.681) is substantially worse, but this reflects sample size reduction (only 797 texts) rather than depth quality effects. Seed variance is minimal ($< 0.002$), confirming reproducibility.

This negative finding is the most robust of our claims: depth fingerprinting is a \textit{diagnostic} tool for provenance characterization, not a curative tool for data curation.

\subsection{Prompt-Type Analysis}
\label{sec:prompt}

Table~\ref{tab:prompt} reports performance by prompt type.

\begin{table}[t]
\centering
\small
\begin{tabular}{lcccc}
\toprule
\textbf{Prompt} & \textbf{CV F1} & \textbf{Ret.} & \textbf{hapax $\rho$} & \textbf{ent. $\rho$} \\
\midrule
Continuation & 0.781 & 0.831 & $-0.74$ & $-0.56$ \\
Rewriting & 0.737 & 0.939 & $-0.80$ & $-0.48$ \\
Elaboration & 0.730 & 0.986 & $-0.83$ & $-0.47$ \\
\bottomrule
\end{tabular}
\caption{Prompt-type comparison (Pythia-2.8B). Ret.\ = depth retention (fraction of chains reaching depth 4). Continuation shows strongest classification signal; elaboration preserves more chains but with weaker signal.}
\label{tab:prompt}
\end{table}

Continuation chains yield the strongest depth signal (CV F1 = 0.781), while elaboration is weakest (0.730)---a 5.1pp gap reflecting moderate prompt-type dependence. The leave-one-prompt-type-out (LOPTO) average F1 of 0.753 confirms reasonable cross-prompt transfer.

\subsection{Paraphrastic Comparison}
\label{sec:paraphrastic}

To connect with attractor cycle theory, we compare continuation and paraphrastic chains (Table~\ref{tab:para}).

\begin{table}[t]
\centering
\small
\begin{tabular}{lccccc}
\toprule
\textbf{Chain} & \textbf{Bin.~F1} & \textbf{3-cl.} & \textbf{hpx $\rho$} & \textbf{ent $\rho$} & \textbf{Conv.} \\
\midrule
Contin. & 0.695 & 0.762 & $-0.63$ & $-0.60$ & $>$d4 \\
Paraph. & 0.572 & 0.580 & $-0.24$ & $-0.12$ & ${\sim}$d2.5 \\
\bottomrule
\end{tabular}
\caption{Paraphrastic vs continuation chains. Conv.\ = approximate attractor convergence depth. Paraphrastic chains converge to attractor states by ${\sim}$depth 2.5, severely attenuating depth signal.}
\label{tab:para}
\end{table}

Paraphrastic chains yield binary F1 of only 0.572 ($\Delta = -12.3$pp vs continuation), with hapax $|\rho|$ dropping from 0.63 to 0.24. The semantic-preservation constraint locks paraphrastic distributions onto an attractor manifold by ${\sim}$depth 2.5, consistent with prior work on iterative paraphrasing. This confirms that our DDR characterization applies specifically to non-paraphrastic chains where the absence of semantic constraints permits sustained monotonic drift.

\subsection{Cross-Model Chains}
\label{sec:cross_model}

To simulate realistic web contamination where different generators contribute at each depth step, we construct cross-model chains (Table~\ref{tab:cross_model}).

\begin{table}[t]
\centering
\small
\begin{tabular}{lccc}
\toprule
\textbf{Evaluation} & \textbf{Bin.~F1} & \textbf{95\% CI} & \textbf{hpx $\rho$} \\
\midrule
Cross-model (gen-LOGO) & 0.830 & [.815, .845] & $-0.68$ \\
Cross-model (comp-LOGO) & 0.764 & [.745, .783] & $-0.65$ \\
Single-model (LOGO) & 0.694 & [.688, .700] & $-0.63$ \\
\bottomrule
\end{tabular}
\caption{Cross-model vs single-model chains. Architecture heterogeneity amplifies depth signal (+13.6pp binary F1), partly attributable to architecture-mixing effects layered on top of depth-intrinsic drift.}
\label{tab:cross_model}
\end{table}

Cross-model chains yield binary F1 = 0.830 under generator-level LOGO, +13.6pp over single-model LOGO. Architecture differences at each depth step introduce additional distributional shifts that amplify the depth signal. The composition-level LOGO (0.764) is lower, indicating that specific generator combinations do not fully transfer.

\subsection{Multi-Domain Analysis}
\label{sec:domains}

Table~\ref{tab:domains} shows binary F1 across four seed domains.

\begin{table}[t]
\centering
\small
\begin{tabular}{lccc}
\toprule
\textbf{Domain} & \textbf{Binary F1} & \textbf{95\% CI} & \textbf{Seeds} \\
\midrule
Wikipedia & 0.815 & [0.798, 0.832] & 1000 \\
Code & 0.754 & [0.737, 0.770] & 500 \\
News & 0.746 & [0.726, 0.764] & 500 \\
Academic & 0.739 & [0.720, 0.757] & 500 \\
\midrule
Cross-domain LOGO & 0.700 & --- & --- \\
\bottomrule
\end{tabular}
\caption{Multi-domain binary F1. Max gap vs Wikipedia is 7.6pp. Cross-domain LOGO (0.700) confirms depth signal is domain-robust.}
\label{tab:domains}
\end{table}

Wikipedia yields the highest F1 (0.815), likely due to distributional alignment with pretraining data. Academic text is lowest (0.739) because its inherently specialized vocabulary compresses the feature margin. The cross-domain LOGO (0.700) is close to the overall LOGO (0.694), confirming moderate domain independence.

\subsection{Finer-Grained Ordinal Classification}
\label{sec:ordinal}

Table~\ref{tab:ordinal} tests whether finer granularity is achievable.

\begin{table}[t]
\centering
\small
\begin{tabular}{lccc}
\toprule
\textbf{Method} & \textbf{5-cl.~F1} & \textbf{3-cl.~F1} & \textbf{$\Delta$} \\
\midrule
XGBoost & 0.500 & 0.762 & $-0.262$ \\
CORAL (ordinal) & 0.500 & --- & --- \\
Frank-Hall (ordinal) & 0.490 & --- & --- \\
\bottomrule
\end{tabular}
\caption{Ordinal granularity comparison. 5-class F1 drops to ${\sim}$0.50 regardless of method. Adjacent-depth confusion rates: gen-1/gen-2 44\%, gen-2/gen-3 53\%, gen-3/gen-4 59\%.}
\label{tab:ordinal}
\end{table}

Five-class classification yields F1 ${\sim}$0.50, and ordinal methods (CORAL, Frank-Hall) provide no improvement over standard XGBoost. Gen-0 is easily identified (F1 = 0.95), but deeper classes are mutually confusable: gen-4 F1 is only 0.28. The bottleneck is feature resolution, not classifier architecture---DDR feature drift slopes flatten after depth 2, making adjacent deep generations indistinguishable. This establishes 3-class (gen-0/gen-1/gen-2+) as the practical resolution ceiling, with principled boundaries aligned to feature discriminability.

\subsection{Scale Effect}
\label{sec:scale}

Adding Qwen3-8B (Table~\ref{tab:scale}) tests whether depth signal persists in larger generators.

\begin{table}[t]
\centering
\small
\begin{tabular}{llcccc}
\toprule
\textbf{Generator} & \textbf{Size} & \textbf{Bin.~F1} & \textbf{hpx $\rho$} & \textbf{ent $\rho$} \\
\midrule
Pythia-1B & 1B & 0.856 & $-0.63$ & $-0.60$ \\
Qwen3-0.6B & 0.6B & 0.824 & $-0.58$ & $-0.56$ \\
Pythia-2.8B & 2.8B & 0.804 & $-0.61$ & $-0.60$ \\
Qwen3-1.7B & 1.7B & 0.805 & $-0.59$ & $-0.55$ \\
OLMo-2-1B & 1.5B & 0.764 & $-0.60$ & $-0.50$ \\
Qwen3-8B & 8B & 0.683 & $-0.56$ & $-0.48$ \\
\bottomrule
\end{tabular}
\caption{Scale effect (6-gen LOGO held-out results). Qwen3-8B shows moderate degradation (0.683) but DDR structure is qualitatively preserved. 6-gen LOGO average: 0.789.}
\label{tab:scale}
\end{table}

Qwen3-8B held-out binary F1 = 0.683---lower than smaller generators but still well above random (0.50). The DDR three-tier structure is qualitatively preserved: hapax ($-0.56$) and TTR ($-0.54$) remain strongly drifting, while tail decay ($-0.09$) stays near-stationary. Surface features are relatively more robust to scale (hapax $|\rho|$ drops 11\%) than reference features (entropy mean $|\rho|$ drops 20\%).

\subsection{Robustness to Simulated Editing}
\label{sec:robustness}

Table~\ref{tab:editing} tests robustness to simulated human editing.

\begin{table}[t]
\centering
\small
\begin{tabular}{lccc}
\toprule
\textbf{Edit Type} & \textbf{Bin.~F1} & \textbf{Drop} & \textbf{95\% CI} \\
\midrule
None (baseline) & 0.695 & --- & [0.679, 0.711] \\
Random 5\% & 0.650 & $-4.5$ & [0.633, 0.667] \\
Random 10\% & 0.614 & $-8.1$ & [0.596, 0.632] \\
Random 20\% & 0.591 & $-10.4$ & [0.573, 0.609] \\
LLM rephrase & 0.647 & $-4.5$ & [0.631, 0.662] \\
\bottomrule
\end{tabular}
\caption{Robustness to simulated editing. The signal degrades gracefully: $-4.5$pp at 5\% random or LLM rephrase, $-10.4$pp at 20\% random. Even at 20\% replacement, F1 remains above random baseline.}
\label{tab:editing}
\end{table}

The depth signal degrades gracefully under simulated editing. LLM-simulated rephrasing (${\sim}$1.7 sentences replaced per document) produces comparable degradation to 5\% random word replacement ($-4.5$pp each). Even at 20\% random replacement (F1 = 0.591), classification remains substantially above chance. Surface features are most vulnerable to word-level perturbation, while reference features are more robust because sentence-level perplexity is less sensitive to individual token replacements.

\subsection{End-to-End Classifier Comparison}
\label{sec:e2e}

Table~\ref{tab:e2e} compares the lightweight 7-feature ensemble against QLoRA fine-tuned end-to-end classifiers.

\begin{table}[t]
\centering
\small
\begin{tabular}{llccc}
\toprule
\textbf{Method} & \textbf{Eval} & \textbf{3-cl.} & \textbf{Bin.} & \textbf{Cost} \\
\midrule
7-feat XGBoost & LOGO & 0.762 & 0.694 & 1$\times$ \\
QLoRA Qwen-Inst & CV & 0.907 & 0.867 & 12$\times$ \\
QLoRA OLMo-Inst & CV & 0.888 & 0.836 & 12$\times$ \\
QLoRA Qwen-Base & CV & 0.904 & 0.860 & 12$\times$ \\
QLoRA OLMo-Base & CV & 0.859 & 0.793 & 12$\times$ \\
\bottomrule
\end{tabular}
\caption{End-to-end (QLoRA) vs handcrafted features. QLoRA is evaluated under within-generator CV (not LOGO), making direct comparison unfair. QLoRA 5-class F1 ranges 0.540--0.598, confirming the resolution ceiling is feature-intrinsic, not method-limited.}
\label{tab:e2e}
\end{table}

QLoRA classifiers achieve substantially higher within-distribution performance (best 3-class F1 = 0.907), but are evaluated under within-generator CV, not LOGO---a fundamentally easier setting. End-to-end models are expected to degrade substantially under LOGO due to overfitting to generator-specific artifacts. The 7-feature ensemble's advantages are interpretability, efficiency (CPU-only, 1$\times$ cost vs 12$\times$ with 8.5GB GPU), and demonstrated cross-generator generalization under LOGO. Even with QLoRA, 5-class F1 remains 0.540--0.598, confirming the resolution ceiling is feature-intrinsic.

\subsection{Text Length Confound Analysis}
\label{sec:length}

\begin{table}[t]
\centering
\small
\begin{tabular}{lll}
\toprule
\textbf{Analysis} & \textbf{Result} & \textbf{Conclusion} \\
\midrule
$\rho$(length, depth) & [$-0.07$, $+0.13$] & Weak ($|\rho| < 0.14$) \\
Length-norm.\ drift & preserved, stronger & Depth $\perp$ length \\
Len-matched F1 & 0.724 vs 0.770 & Sample-size effect \\
Rand.\ subsample & ${\sim}$0.730 (same $n$) & Confirms above \\
\bottomrule
\end{tabular}
\caption{Length confound analysis. Text length is not a confound: depth-length correlation is weak, and length normalization strengthens feature drift.}
\label{tab:length}
\end{table}

Text length is not a confound for depth classification (Table~\ref{tab:length}). Depth-length Spearman correlations are weak ($|\rho| < 0.14$). After length normalization, TTR and hapax depth drift are preserved and strengthened. The 4.6pp F1 drop under length-matched subsampling is attributable to reduced sample size (random subsample control at the same $n$ yields ${\sim}$0.730), not confound removal.

\subsection{LD-Score as Additional Feature}
\label{sec:ldscore}

\begin{table}[t]
\centering
\small
\begin{tabular}{lcc}
\toprule
\textbf{Feature Set} & \textbf{LOGO Bin.~F1} & \textbf{$\Delta$} \\
\midrule
7 features & 0.694 & --- \\
7 features + LD-Score & 0.700 & $+0.6$ \\
LD-Score alone & 0.721 & --- \\
\bottomrule
\end{tabular}
\caption{LD-Score as additional feature. Marginal gain (+0.6pp) reflects high redundancy with hapax ratio ($r = 0.75$). LD-Score alone achieves 0.721 LOGO F1 but with wide cross-generator variance [0.588, 0.770].}
\label{tab:ldscore}
\end{table}

Adding LD-Score (a character-level distribution signature) as an 8th feature provides only +0.6pp marginal gain (Table~\ref{tab:ldscore}), reflecting high redundancy with hapax ratio ($r = 0.75$). LD-Score's strong ordinal signal ($\rho = -0.79$) confirms that depth-induced distributional drift manifests at the character level, but its information is largely captured by existing surface features. We retain the 7-feature design for parsimony.

%------------------------------------------------------------
\section{Conclusion}
\label{sec:conclusion}
%------------------------------------------------------------

We formulated generation depth classification as a task beyond binary AI-text detection and characterized differential depth response (DDR) across 5 base-model generators from 3 families. A lightweight 7-feature ensemble achieves leak-free binary F1 = 0.694 under LOGO evaluation, ranging 0.656--0.724 per generator, and 0.830 on cross-model chains. DDR is robust across tested decoding strategies (greedy, top-k) and three reference models from different families.

Our most robust finding is negative: depth-aware data curation is statistically equivalent to random subsampling for downstream CLM quality. Depth fingerprinting is diagnostic---useful for provenance characterization---not curative for data selection.

The signal degrades substantially under instruction tuning (OLMo-Instruct LOGO F1 = 0.338 vs 0.763 base), restricting current scope to base models. This degradation is itself a finding: instruction tuning compresses the distributional variation DDR exploits. Future work should investigate whether alternative feature families can recover depth signal in instruct models, extend evaluation to larger frontier models ($>$14B), and incorporate additional decoding strategies (nucleus sampling, beam search) to strengthen robustness claims.

%------------------------------------------------------------
\section*{Limitations}
%------------------------------------------------------------

Our study has several limitations. First, all experiments use base models; instruction-tuned models show severely degraded signal, and extending depth fingerprinting to the instruct regime remains open. Second, our decoding robustness analysis covers greedy and top-k sampling; claims should be considered preliminary until validated on nucleus sampling and beam search. Third, the downstream utility experiment uses only perplexity on a single CLM task with 3 seeds---additional tasks and metrics are needed to confirm the ``diagnostic not curative'' conclusion broadly.

\bibliography{custom}

\end{document}
